% GENERATED FILE. Edit canonical lessons, then run npm run generate:books.
% canonical-source-hash: fnv1a64:ef1981de40e21218
% canonical-lessons: PA-C99-adhirat, PA-C99-tarka, PA-C99-parso, PA-C99-garmi, PA-C99-sardi

\chapter{Midnight, Daybreak, Day after tomorrow, Summer, Winter}
\label{ch:pa-midnight-daybreak-day-after-tomorrow-summer-winter}
\hlchaptermodality{99}

\begin{chapteropening}
\textbf{By the end of this chapter:} \emph{I can say midnight, daybreak, the day after tomorrow, summer and winter.}

Complete the last lesson of chapter 99: I can say midnight, daybreak, the day after tomorrow, summer and winter.
\end{chapteropening}

\section[\texorpdfstring{\pa{ਅੱਧੀ} \pa{ਰਾਤ} (adhī rāt)}{ਅੱਧੀ ਰਾਤ (adhī rāt)}]{\pa{ਅੱਧੀ} \pa{ਰਾਤ} (adhī rāt) — midnight}

\label{lesson:PA-C99-adhirat}



\begin{quote}
Before the new one: say the Punjabi for a way, a route, then the Punjabi for a country.

Say \textquotedblleft{}I live in Amritsar.\textquotedblright{} (\emph{Maiṁ Amritsar rahindā hāṁ}: the place-name needs no \textquotedblleft{}in,\textquotedblright{} and the verb comes last.)
\end{quote}

\subsection*{You'll want to know: \pa{ਅੱਧੀ} \pa{ਰਾਤ}}
\textbf{\pa{ਅੱਧੀ} \pa{ਰਾਤ}} — \emph{adhī rāt} — \textquotedblleft{}midnight\textquotedblright{}.

\begin{cousinweb}
Both words are inherited. \emph{Adhī}, \textquotedblleft{}half\textquotedblright{}, goes with Punjabi \emph{addhā}, which CDIAL lists under Sanskrit \emph{ardha}, \textquotedblleft{}half\textquotedblright{} (CDIAL 644), and \emph{rāt} comes from \emph{rātrī}, \textquotedblleft{}night\textquotedblright{} (CDIAL 10702).
\end{cousinweb}

\begin{grammarlens}[title={What you've built}]
One more word for this chapter.
\end{grammarlens}

\subsection*{Your turn}
\begin{itemize}
\raggedright
  \item \emph{Say it:} \emph{adhī rāt}
  \item \emph{Say it:} \emph{adhī rāt}, once more
  \item \emph{Say it:} say \emph{adhī rāt}
  \item \emph{Recall:} say the Punjabi for a way, a route, then the Punjabi for a country, then say \emph{adhī rāt} again
\end{itemize}

\begin{tcolorbox}[breakable,colback=teal!4,colframe=teal!35!black,title={Before you move on}]
What does \pa{ਅੱਧੀ} \pa{ਰਾਤ} mean? (\textquotedblleft{}Midnight\textquotedblright{}.) Say it once more.
\end{tcolorbox}
\section[\texorpdfstring{\pa{ਤੜਕਾ} (taṛkā)}{ਤੜਕਾ (taṛkā)}]{\pa{ਤੜਕਾ} (taṛkā) — daybreak, early morning}

\label{lesson:PA-C99-tarka}



\begin{quote}
Before the new one: say the Punjabi for a country, then the Punjabi for midnight.
\end{quote}

\subsection*{You'll want to know: \pa{ਤੜਕਾ}}
\textbf{\pa{ਤੜਕਾ}} — \emph{taṛkā} — \textquotedblleft{}daybreak, early morning\textquotedblright{}.

\begin{cousinweb}
\textbf{\pa{ਤੜਕਾ}} (\emph{taṛkā}) matches Hindi \emph{taṛkā}, \textquotedblleft{}dawn\textquotedblright{}, which CDIAL lists under the imitative Sanskrit \emph{traṭ}, \textquotedblleft{}crackle\textquotedblright{}, comparing the English \textquotedblleft{}crack of dawn\textquotedblright{} (CDIAL 5988). The same entry gives Punjabi \emph{taṛakṇā}, \textquotedblleft{}to crack, crackle\textquotedblright{}.
\end{cousinweb}

\begin{grammarlens}[title={What you've built}]
One more word for this chapter.
\end{grammarlens}

\subsection*{Your turn}
\begin{itemize}
\raggedright
  \item \emph{Say it:} \emph{taṛkā}
  \item \emph{Say it:} \emph{taṛkā}, once more
  \item \emph{Say it:} say \emph{taṛkā}
  \item \emph{Recall:} say the Punjabi for a country, then the Punjabi for midnight, then say \emph{taṛkā} again
\end{itemize}

\begin{tcolorbox}[breakable,colback=teal!4,colframe=teal!35!black,title={Before you move on}]
What does \pa{ਤੜਕਾ} mean? (\textquotedblleft{}Daybreak, early morning\textquotedblright{}.) Say it once more.
\end{tcolorbox}
\section[\texorpdfstring{\pa{ਪਰਸੋਂ} (parsõ)}{ਪਰਸੋਂ (parsõ)}]{\pa{ਪਰਸੋਂ} (parsõ) — the day after tomorrow}

\label{lesson:PA-C99-parso}



\begin{quote}
Before the new one: say the Punjabi for midnight, then the Punjabi for daybreak.
\end{quote}

\subsection*{You'll want to know: \pa{ਪਰਸੋਂ}}
\textbf{\pa{ਪਰਸੋਂ}} — \emph{parsõ} — \textquotedblleft{}the day after tomorrow\textquotedblright{}.

\begin{cousinweb}
\textbf{\pa{ਪਰਸੋਂ}} (\emph{parsõ}) is a tadbhava, a word inherited from Sanskrit and reshaped by sound change. It comes from \emph{paraśvas}, \textquotedblleft{}the day after tomorrow\textquotedblright{} (CDIAL 7799j). Punjabi \emph{parsõ} means both the day after tomorrow and the day before yesterday.
\end{cousinweb}

\begin{grammarlens}[title={What you've built}]
One more word for this chapter.
\end{grammarlens}

\subsection*{Your turn}
\begin{itemize}
\raggedright
  \item \emph{Say it:} \emph{parsõ}
  \item \emph{Say it:} \emph{parsõ}, once more
  \item \emph{Say it:} say \emph{parsõ}
  \item \emph{Recall:} say the Punjabi for midnight, then the Punjabi for daybreak, then say \emph{parsõ} again
\end{itemize}

\begin{tcolorbox}[breakable,colback=teal!4,colframe=teal!35!black,title={Before you move on}]
What does \pa{ਪਰਸੋਂ} mean? (\textquotedblleft{}The day after tomorrow\textquotedblright{}.) Say it once more.
\end{tcolorbox}
\section[\texorpdfstring{\pa{ਗਰਮੀ} (garmī)}{ਗਰਮੀ (garmī)}]{\pa{ਗਰਮੀ} (garmī) — summer, the heat}

\label{lesson:PA-C99-garmi}



\begin{quote}
Before the new one: say the Punjabi for daybreak, then the Punjabi for the day after tomorrow.
\end{quote}

\subsection*{You'll want to know: \pa{ਗਰਮੀ}}
\textbf{\pa{ਗਰਮੀ}} — \emph{garmī} — \textquotedblleft{}summer, the heat\textquotedblright{}.

\begin{cousinweb}
\textbf{\pa{ਗਰਮੀ}} (\emph{garmī}) is a loan, usually traced to Persian \emph{garmī}, \textquotedblleft{}heat\textquotedblright{}. It is Persian \emph{garm}, \textquotedblleft{}hot\textquotedblright{}, with the ending \emph{-ī}.
\end{cousinweb}

\begin{grammarlens}[title={What you've built}]
One more word for this chapter.
\end{grammarlens}

\subsection*{Your turn}
\begin{itemize}
\raggedright
  \item \emph{Say it:} \emph{garmī}
  \item \emph{Say it:} \emph{garmī}, once more
  \item \emph{Say it:} say \emph{garmī}
  \item \emph{Recall:} say the Punjabi for daybreak, then the Punjabi for the day after tomorrow, then say \emph{garmī} again
\end{itemize}

\begin{tcolorbox}[breakable,colback=teal!4,colframe=teal!35!black,title={Before you move on}]
What does \pa{ਗਰਮੀ} mean? (\textquotedblleft{}Summer, the heat\textquotedblright{}.) Say it once more.
\end{tcolorbox}
\section[\texorpdfstring{\pa{ਸਰਦੀ} (sardī)}{ਸਰਦੀ (sardī)}]{\pa{ਸਰਦੀ} (sardī) — winter, the cold}

\label{lesson:PA-C99-sardi}



\begin{quote}
Before the new one: say the Punjabi for the day after tomorrow, then the Punjabi for summer.
\end{quote}

\subsection*{You'll want to know: \pa{ਸਰਦੀ}}
\textbf{\pa{ਸਰਦੀ}} — \emph{sardī} — \textquotedblleft{}winter, the cold\textquotedblright{}.

\begin{cousinweb}
\textbf{\pa{ਸਰਦੀ}} (\emph{sardī}) is a loan, usually traced to Persian \emph{sardī}, \textquotedblleft{}cold\textquotedblright{}. It is Persian \emph{sard}, \textquotedblleft{}cold\textquotedblright{}, with the ending \emph{-ī}.
\end{cousinweb}

\begin{grammarlens}[title={What you've built}]
One more word for this chapter.
\end{grammarlens}

\subsection*{Your turn}
\begin{itemize}
\raggedright
  \item \emph{Say it:} \emph{sardī}
  \item \emph{Say it:} \emph{sardī}, once more
  \item \emph{Say it:} say \emph{sardī}
  \item \emph{Recall:} say the Punjabi for the day after tomorrow, then the Punjabi for summer, then say \emph{sardī} again
\end{itemize}

\begin{tcolorbox}[breakable,colback=teal!4,colframe=teal!35!black,title={Before you move on}]
What does \pa{ਸਰਦੀ} mean? (\textquotedblleft{}Winter, the cold\textquotedblright{}.) Say it once more.
\end{tcolorbox}
